\documentclass[10pt]{article}

\usepackage[utf8]{inputenc}

\usepackage[T1]{fontenc}

\usepackage{amsmath}

\usepackage{amsfonts}

\usepackage{amssymb}

\usepackage[version=4]{mhchem}

\usepackage{stmaryrd}

\usepackage{bbold}



\begin{document}

нейное преобразование пространства $E$, представляющееся в базисе $\left\{e_{i}\right\}$ матрицей вида



\[

\left\|\begin{array}{cc}

I_{m} & 0 \\

0 & \beta I_{m}

\end{array}\right\|

\]



( $\beta$ - ненулевой әлемент поля $K)$.



C. г. $\mathrm{Sp}_{2 m}(K)$ совпадает с группой $K$-точек линейной алгебраич. группы $\mathrm{Sp}_{2 m}$, задаваемой уравнением $A^{\perp} J_{m} A=J_{m}$. Эта алгебраич. группа, тоже называемая С. Г., является простой односвязной линейной алгебраич. группой типа $C_{m}$, ее размерность равна $2 m^{2}+m$.



В случае, когда $K=\mathbb{C}$ или $\mathbb{R}$, С. г. $\operatorname{Sp}_{2 m}(K)$ есть связная односвязная простая комплексная (соответственно вещественная) группа Ли. Группа $\mathrm{Sp}_{2 m}(\mathbb{R})$ является одной из вещественных форм комплекснрй C. г. $\mathrm{Sp}_{2 m}(\mathbb{C})$. Остальные вещественные формы әтой группы тоже иногда называют С. г. Это - подгруппы $\mathrm{Sp}(p, q), p, q \geqslant 0, p+q=m$, выделяемые из групшы $\mathrm{Sp}_{2 m}(\mathbb{C})$ условием сохранения эрмитовой формы вида



\[

\sum_{i=1}^{2 m} \varepsilon_{i} \bar{z}_{i} \bar{z}_{i}

\]



где $\varepsilon_{i}=1$ при $1 \leqslant i \leqslant p$ и $m+1 \leqslant i \leqslant m+p$ п $\varepsilon_{i}=-1$ при остальных $i$. Группа $\mathrm{Sp}(0, m)$ - компактная вещественная форма комплексной С. г. $\mathrm{Sp}_{2 m}(\mathbb{C})$. С. г. $\operatorname{Sp}(p, q)$ изоморфна группе всех линейных преобразований правого векторного пространства $\mathrm{H}^{m}$ над телом кватернионов $H$ размерности $m=p+q$, к-рые сохраняют кватернионную эрмитову форму индекса $\min (p, q)$, то есть форму вида



\[

(x, y)=\sum_{i=1}^{p} x_{i} \bar{y}_{i}-\sum_{i=p+1}^{m} x_{i} \bar{y}_{i},

\]



где $x=\left(x_{1}, x_{2}, \ldots, x_{m}\right), y=\left(y_{1}, y_{2}, \ldots, y_{m}\right) \in \mathbb{H}^{m}$, а черта означает переход к сопряженному кватерниону.



Лит.: [1] А р т и н Э., Геометрическая алгебра, пер. с англ., М., 1969; [2] Б у р б а к и H., Алгебра. Модули, кольца, формы, пер. с франц., М., 1966; [3] Д ь ё д о н н е ЖЖ., Геометрия классических групп, пер. с франц., М., 1974; [4] Х е л г а с о в С., Диффференциальная геометрия и симметрические пространства, пер. с англ., М., 1964; [5] III е в а л л е К., Теория групп Ли, пер. с англ., т. 1, М., 1948.



СИМПЛЕКТИЧЕСКАЯ СВЯЗНОСТБ - афбинная связность на гладком многообразии $M$ размерности $2 n$, обладающая ковариантно постоянной относительно нее невырожденной 2-формой Ф. Если афффинная связность на $M$ задана с помощью локальных форм связности



\[

\begin{aligned}

& \omega^{i}=\Gamma_{k}^{i} d k^{k}, \quad \operatorname{det}\left\|\Gamma_{k}^{i}\right\| \neq 0, \\

& \omega_{j}^{i}=\Gamma_{j k}^{i} \omega^{k}

\end{aligned}

\]



и



\[

\Phi=a_{i j} \omega^{i} \wedge \omega^{j}, a_{i j}=-a_{j i}

\]



то условие ковариантного постоянства Ф выражается в виде



\[

d a_{i j}=a_{k j} \omega_{i}^{k}+a_{i k} \omega_{j}^{k} .

\]



2-форма Ф определяет на $M$ симплектич. структуру (или почти гамильтонову структуру), превращая каждое касательное пространство $T_{x}(M)$ в симплектич. пространство с кососимметрич. скалярным произведением $\Phi(X, Y)$. С. с. можно определить так же, как аффинную связность на $M$, сохраняющую при параллельном перенесении векторов это произведение. Репер в каждом $T_{x}(M)$ может быть выбран так, что



\[

\Phi=2 \sum_{\alpha=1}^{n} \omega^{\alpha} \wedge \omega^{n+\alpha} .

\]



Все такие реперы образуют главное расслоенное пространство над $M$, структурной группой к-рого является симплектич. группа. С. с.-әто связность в әтом главном расслоенном пространстве. Существуют многообразия $M$ qетной размерности, на к-рых нет невырожденной глобально определенной 2-формы $\Phi$ и, следовательно, нет С. с. Однако если Ф существует, то С. с., относительно к-рой Ф ковариантно постоянно, определяется неоднозначно.



Лит.: [1] С т е р п б е р г С., Лекцции по дифференциальной геометрии, пер. с англ., М., 1970 . Ю. Г. Лулисте. тезимальная структура 1-го порядка на четномерном гладком ориентируемом многообразии $M^{2 n}$, к-рая оІІределяется заданием на $M^{2 n}$ невырожденной 2-формы $Ф$. В каждом касательном пространстве $T_{x}\left(M^{2 n}\right)$ возникает структура симплектич. пространства с кососимметрическим скалярным произведением $\Phi(X, Y)$. Bсе касательные к $M^{2 n}$ реперы, адаптированные к С. с. (т. е. реперы, относительно к-рых Ф имеет канонич. вид



\[

\left.\Phi=2 \sum_{\alpha=1}^{n} \omega^{\alpha} \wedge \omega^{n+\alpha}\right),

\]



образуют главное расслоенное пространство над $M^{2 n}$, структурной группой к-рого является симплектич. группа $\operatorname{Sp}(n)$. Вообще, задание С. с. на $M^{2 n}$ равносильно заданию $\operatorname{Sp}(n)$-структуры на $M^{2 n}$, как нек-рой G-структуры.



На $M^{2 n}$ со С. с. существует изоморфизм между модулями векторных полей и 1-форм на $M^{2 n}$, к-рый векторному полю $X$ ставит в соответствие 1-форму $\omega_{X}: Y \rightarrow \Phi(X, Y)$. Образ скобки Ли $[X, Y]$ наз. при этом скобкой Пуассона $\left[\omega_{X}, \omega_{Y}\right]$; в частности, когда $\omega_{X}$ и $\omega_{Y}$ полные дифференциалы, получается понятие скобки Пуассона двух функций на $M^{2 n}$, к-рое обобщает соответствующее классич. понятие.



C. с. наз. п о ч т и г а м и л ь т о но в о й с т р у кт у р о й, а если Ф замкнута, т. е. $d \Phi=0$, то г а м и л ьт о н о во й с т р укт у ро й; впрочем, иногда условие $d \Phi=0$ включают в определение С. с. Эти структуры, находящие применения в глобальной аналитич. механике, основаны на том факте, что на касательном расслоенном пространстве $T^{*}(M)$ любого гладкого многообразия $M$ существует каноническая гамильтонова структура. Она определяется формой $\Phi=d \theta$, где 1-форма $\theta$ на $T^{*}(M)$, наз. фо р м о й Ј и ув и л л я, задается следующим образом: $\theta_{t}\left(X_{t}\right)=$ $=u\left(\pi_{*} X_{u}\right)$ для любого касательного вектора $X_{u}$ в точке $u \in T^{*}(M)$, где $\pi-$ проекция $T^{*}(M) \rightarrow M$. Если на $M$ выбраны локальные координаты $x^{1}, \ldots, x^{n}$ и $u=y_{i}(u) d x_{\pi(u)}^{i}$, то $\theta=y_{i} d x^{i}$, вследствие чего $\dot{\Phi}=d y_{i} \wedge d x^{i}$. В классич. механике $M$ интерпретируется как конфигурационное пространство, а $T^{*}(M)$ как фазовое пространство.



Векторное поле $X$ на $M^{2 n}$ с гамильтоновой структурой наз. г а м и л ь тон о в ы м (или г а м и л ь тонов о й с и с темо й), если 1-форма $\omega_{X}$ замкнута. Если она, кроме того, точна, т. е. $\omega_{X}=-d H$, то функция $H$ на $M^{2 n}$ наз. г а м и л ь т о н и а но м и является обобщением соответствующего классического понятия.



Лит.: [1] С т е р в 6 е p г С., Лекции по дифференциальной геометрии, пер. с англ., М., $1970 ;[2] \Gamma$ о д б и й о н К., Диффранияаль геометрия и аналитическая механика, пер. с



СИМПЛЕКТИЧЕСКОЕ МНОГООБРАЗИЕ - Многообразие, снабженное симплектической структурой.



СИМПЛЕКТИЧЕСКОЕ ПРОСТРАНСТВО - нечетномерное проективное пространство $P_{2 n+1}$ над полем $k$ с заданной в нем инволюционной корреляцией - нульсистемой; обовначается $\mathrm{Sp}_{2 n+1}$.



Пусть характеристика поля $h$ не равна 2. Абсолютная нуль-система в $\mathrm{Sp}_{2 n+1}$ всегда может быть записана в виде $u_{i}=a_{i j} x^{j}$, где $\left\|a_{i j}\right\|-$ кососимметрич. матрица





\end{document}